\documentclass[aspectratio=169,9pt]{beamer}

\usetheme{default}
\usefonttheme{professionalfonts}
\setbeamertemplate{navigation symbols}{}
\setbeamertemplate{footline}{\hfill\scriptsize\insertframenumber/3\hspace{2mm}\vspace{1mm}}

\usepackage{amsmath,amssymb,bm,mathtools}
\usepackage{xcolor}
\usepackage[most]{tcolorbox}

\definecolor{deepblue}{RGB}{0,70,150}
\definecolor{softblue}{RGB}{232,242,255}
\definecolor{softgreen}{RGB}{235,250,238}
\definecolor{softorange}{RGB}{255,245,230}
\definecolor{softpurple}{RGB}{245,238,255}
\definecolor{darkred}{RGB}{160,35,25}
\definecolor{darkgreen}{RGB}{20,120,60}
\definecolor{darkpurple}{RGB}{95,45,150}

\setbeamerfont{frametitle}{series=\bfseries,size=\Large}
\setbeamercolor{frametitle}{fg=black}
\setbeamertemplate{itemize item}{\color{deepblue}$\blacktriangleright$}
\setlength{\abovedisplayskip}{2pt}
\setlength{\belowdisplayskip}{2pt}
\setlength{\jot}{1pt}

\newtcolorbox{bluebox}[1][]{colback=softblue,colframe=deepblue,boxrule=0.75pt,arc=1.4mm,left=1.3mm,right=1.3mm,top=0.8mm,bottom=0.8mm,#1}
\newtcolorbox{greenbox}[1][]{colback=softgreen,colframe=darkgreen,boxrule=0.75pt,arc=1.4mm,left=1.3mm,right=1.3mm,top=0.8mm,bottom=0.8mm,#1}
\newtcolorbox{orangebox}[1][]{colback=softorange,colframe=darkred,boxrule=0.75pt,arc=1.4mm,left=1.3mm,right=1.3mm,top=0.8mm,bottom=0.8mm,#1}
\newtcolorbox{purplebox}[1][]{colback=softpurple,colframe=darkpurple,boxrule=0.75pt,arc=1.4mm,left=1.3mm,right=1.3mm,top=0.8mm,bottom=0.8mm,#1}

\newcommand{\Uvec}{\bm{U}}
\newcommand{\Fvec}{\bm{F}}
\newcommand{\Cvec}{\bm{C}}
\newcommand{\Evec}{\bm{E}}
\newcommand{\Jvec}{\bm{J}}
\newcommand{\Jn}{\bm{J}_n}
\newcommand{\Jp}{\bm{J}_p}
\newcommand{\qrad}{\dot{q}_{\mathrm{rad}}}
\newcommand{\qth}{\dot{q}_{\mathrm{th}}}
\newcommand{\grad}{\nabla}
\newcommand{\diverg}{\nabla\!\cdot}
\newcommand{\epss}{\varepsilon_s}
\newcommand{\rhom}{\rho_m}
\newcommand{\cp}{c_p}

\begin{document}

% -----------------------------------------------------------------------------
\begin{frame}{Module 6: one coupled state, four interacting maps}
\footnotesize

\begin{bluebox}
\centering\small
\textbf{Goal:} solve defect concentration, electrostatic potential, temperature, and carrier transport as one self-consistent device problem.
\end{bluebox}

\vspace{1.2mm}
\begin{columns}[T,totalwidth=\textwidth]
\begin{column}{0.46\textwidth}
\textbf{Coupled state vector}
\begin{equation*}
\boxed{
\Uvec(x,y,t)=
\begin{bmatrix}
C_1,\ldots,C_N,\phi,T,n,p
\end{bmatrix}^{T}}
\end{equation*}
\small
Here, \(C_j\) are defect concentrations, \(\phi\) is electrostatic potential, \(T\) is temperature, and \(n,p\) are carrier densities.

\vspace{1.5mm}
\begin{purplebox}
\textbf{Core idea:} each map is both an output and an input. Defects, fields, heat, and carriers change the coefficients and source terms of one another.
\end{purplebox}
\end{column}

\begin{column}{0.52\textwidth}
\textbf{Physics blocks embedded in Module 6}
\begin{bluebox}
\scriptsize
\begin{align*}
\frac{\partial C_j}{\partial t}
&=\diverg(D_j\grad C_j)-\diverg(\mu_j z_j q C_j\Evec)+R_j+S_j,\\
\diverg(\epss\grad\phi)&=-\rho,\qquad \Evec=-\grad\phi,\\
\rhom\cp\frac{\partial T}{\partial t}
&=\diverg(k\grad T)+Q_{\mathrm{tot}},\\
\frac{\partial n}{\partial t}
&=\frac{1}{q}\diverg\Jn+G-R,\qquad
\frac{\partial p}{\partial t}=-\frac{1}{q}\diverg\Jp+G-R.
\end{align*}
\end{bluebox}
\small
These are Modules 2--5 written on the same 2D domain with shared fields and consistent source terms.
\end{column}
\end{columns}

\vspace{1.2mm}
\begin{orangebox}
\centering\scriptsize
\textbf{Module 6 is the coupling layer:} it prevents stale fields, inconsistent signs, and one-way assumptions from contaminating the final device prediction.
\end{orangebox}
\end{frame}

% -----------------------------------------------------------------------------
\begin{frame}{Module 6: feedback loops and coupling terms}
\footnotesize

\begin{bluebox}
\centering\small
Radiation damage enters through source terms, but degradation emerges from feedback among space charge, heat, mobility, recombination, and current.
\end{bluebox}

\vspace{1.2mm}
\begin{columns}[T,totalwidth=\textwidth]
\begin{column}{0.48\textwidth}
\textbf{Defects and carriers perturb electrostatics}
\begin{bluebox}
\scriptsize
\[
\rho=q\left[p-n+N_D^+-N_A^-+
\sum_{j=1}^{N}z_jC_j\right]
\]
\[
\diverg(\epss\grad\phi)=-\rho,
\qquad \Evec=-\grad\phi
\]
\end{bluebox}
\small
Charged defects enter the same charge-density budget as dopants and mobile carriers, so a defect cloud bends \(\phi\) and \(\Evec\).

\vspace{1mm}
\textbf{Temperature changes kinetic coefficients}
\begin{orangebox}
\scriptsize
\begin{align*}
D_j(T)&=D_{j0}\exp\!\left(-\frac{E_{D,j}}{k_BT}\right),\\[-1mm]
k_{\mathrm{ann},j}(T)&=k_{0,j}\exp\!\left(-\frac{E_{A,j}}{k_BT}\right).
\end{align*}
\end{orangebox}
\end{column}

\begin{column}{0.50\textwidth}
\textbf{Defects modify transport and recombination}
\begin{greenbox}
\scriptsize
\[
\mu_n(T,\Cvec)=\frac{\mu_{n0}(T)}{1+\sum_j\alpha_{n,j}C_j},
\qquad
\mu_p(T,\Cvec)=\frac{\mu_{p0}(T)}{1+\sum_j\alpha_{p,j}C_j}
\]
\[
N_t(x,y,t)=\sum_j\eta_j C_j(x,y,t)
\]
\[
R_{\mathrm{SRH}}=
\frac{np-n_i^2}{\tau_p(n+n_1)+\tau_n(p+p_1)}
\]
\end{greenbox}

\vspace{0.5mm}
\textbf{Carrier transport feeds heat}
\begin{purplebox}
\scriptsize
\[
Q_{\mathrm{tot}}=\qth+\Jvec\!\cdot\!\Evec+Q_{\mathrm{rec}}+Q_{\mathrm{def}}+Q_{\mathrm{ext}}
\]
\end{purplebox}
\end{column}
\end{columns}

\vspace{1.1mm}
\begin{greenbox}
\centering\scriptsize
\(\qrad \rightarrow C_j,T,n,p \rightarrow \rho,\phi,\Evec \rightarrow \Jn,\Jp,R,Q_{\mathrm{tot}} \rightarrow T,C_j,n,p\). This closed loop is the physical reason Module 6 exists.
\end{greenbox}
\end{frame}

% -----------------------------------------------------------------------------
\begin{frame}{Module 6: FEM residuals, staged coupling, and outputs}
\scriptsize

\begin{bluebox}
\centering
The computational problem is a coupled PDE/DAE-like residual system: time-dependent defect, thermal, and carrier equations constrained by instantaneous Poisson electrostatics.
\end{bluebox}

\vspace{1mm}
\begin{columns}[T,totalwidth=\textwidth]
\begin{column}{0.54\textwidth}
\textbf{Monolithic residual form}
\begin{equation*}
\boxed{
\Fvec\!\left(\Uvec,\frac{\partial\Uvec}{\partial t};\qrad,\mathrm{BCs},\mathrm{params}\right)=\bm{0}}
\end{equation*}
\begin{align*}
F_{C_j}&=\partial_t C_j-\diverg(D_j\grad C_j)+\diverg(\mu_jz_jqC_j\Evec)-R_j-S_j,\\
F_{\phi}&=\diverg(\epss\grad\phi)+\rho,\\
F_T&=\rhom\cp\partial_tT-\diverg(k\grad T)-Q_{\mathrm{tot}},\\
F_n&=\partial_t n-\frac{1}{q}\diverg\Jn-G+R,\\
F_p&=\partial_t p+\frac{1}{q}\diverg\Jp-G+R.
\end{align*}
\begin{orangebox}
\textbf{FEM role:} convert each residual to weak form; carry \(C_j,\phi,T,n,p\) on the same mesh or compatible interpolated meshes.
\end{orangebox}
\end{column}

\begin{column}{0.44\textwidth}
\textbf{First implementation: staggered fixed-point solve}
\begin{enumerate}\itemsep0.4mm
\item update \(C_j^{m+1}\) using \(T^m,\Evec^m,n^m,p^m\),
\item update \(\phi^{m+1}\) from \(\rho(C_j,n,p)\),
\item update \(T^{m+1}\) from \(Q_{\mathrm{tot}}\),
\item update \(n^{m+1},p^{m+1}\) from \(\Evec,T,C_j\),
\item iterate until field changes are small.
\end{enumerate}
\[
\boxed{\frac{\|\Uvec^{(r+1)}-\Uvec^{(r)}\|}{\|\Uvec^{(r)}\|+\epsilon}<\mathrm{tol}}
\]

\vspace{0.5mm}
\textbf{Device-level outputs}
\[
\boxed{I_{\mathrm{term}}=\int_{\Gamma_{\mathrm{contact}}}(\Jn+\Jp)\cdot\bm{n}\,d\Gamma}
\]
\[
\boxed{y_{\mathrm{device}}=\mathcal{G}(I,V_{\mathrm{th}},Q_{\mathrm{coll}},I_{\mathrm{leak}},\ldots)}
\]
These observables feed Module 8 scoring and dominant-path reduction.
\end{column}
\end{columns}
\end{frame}

\end{document}
